% ============================================================
%  课程笔记封面模板 —— 「雾海中的漫游者」风格
%  取材：弗里德里希《雾海上的漫游者》(Wanderer above the Sea of Fog, 1818)：
%        冷灰蓝的雾海、层叠的岩峰剪影、立于绝顶的背影与手杖。
%  与 cover_Impression风.tex 同构（眉题 / 主标题 / 菱形饰线 /
%  副标题 / 底部作者日期），直接 \input{} 复用：
%      % ============================================================
%  课程笔记封面模板 —— 「雾海中的漫游者」风格
%  取材：弗里德里希《雾海上的漫游者》(Wanderer above the Sea of Fog, 1818)：
%        冷灰蓝的雾海、层叠的岩峰剪影、立于绝顶的背影与手杖。
%  与 cover_Impression风.tex 同构（眉题 / 主标题 / 菱形饰线 /
%  副标题 / 底部作者日期），直接 \input{} 复用：
%      % ============================================================
%  课程笔记封面模板 —— 「雾海中的漫游者」风格
%  取材：弗里德里希《雾海上的漫游者》(Wanderer above the Sea of Fog, 1818)：
%        冷灰蓝的雾海、层叠的岩峰剪影、立于绝顶的背影与手杖。
%  与 cover_Impression风.tex 同构（眉题 / 主标题 / 菱形饰线 /
%  副标题 / 底部作者日期），直接 \input{} 复用：
%      % ============================================================
%  课程笔记封面模板 —— 「雾海中的漫游者」风格
%  取材：弗里德里希《雾海上的漫游者》(Wanderer above the Sea of Fog, 1818)：
%        冷灰蓝的雾海、层叠的岩峰剪影、立于绝顶的背影与手杖。
%  与 cover_Impression风.tex 同构（眉题 / 主标题 / 菱形饰线 /
%  副标题 / 底部作者日期），直接 \input{} 复用：
%      \input{cover_雾海漫游者风}
%
%  配色：wf* 前缀（弗里德里希取色）。本文件用 \providecolor 兜底，
%        单独复制到任何笔记本也能编译；想整体换调改这里即可。
%
%  注意：整幅画面绘于 \begin{scope}[shift={(current page.south west)}] 内，
%        所有坐标以「页面左下角为原点、单位 cm」书写；脱离该 scope
%        直接写 (x,y)，图形会落到页面外而不可见。
%
%  可用字段（在主文件导言区用 \newcommand 预定义即可覆盖缺省值）：
%      \notename     主标题（必填，主模板已定义）
%      \noteauthor   作者（必填，主模板已定义）
%      \notecourseen 眉题英文（缺省 Course Notes）
%      \notesubtitle 副标题（缺省「学习笔记」）
%      \noteextra    底部附加信息（缺省不显示）
% ============================================================

% ---- 弗里德里希《雾海上的漫游者》取色（冷灰蓝 + 墨黑剪影） ----
\providecolor{wfSky}{RGB}{188,198,210}      % 高空：灰蓝（上）
\providecolor{wfSky2}{RGB}{214,220,228}     % 近雾：浅灰蓝
\providecolor{wfFog}{RGB}{240,242,244}      % 雾海：白
\providecolor{wfFar}{RGB}{146,158,172}      % 远峰：淡灰蓝
\providecolor{wfRock2}{RGB}{104,106,112}    % 中景岩：中灰
\providecolor{wfRock}{RGB}{58,60,66}        % 前景岩：深灰
\providecolor{wfFigure}{RGB}{28,30,34}      % 人物：墨黑
\providecolor{wfInk}{RGB}{36,42,50}         % 标题：深蓝灰

% ---- 缺省字段（主文件已定义的同名命令不会被覆盖） ----
\providecommand{\notecourseen}{Course Notes}
\providecommand{\notesubtitle}{学\ 习\ 笔\ 记}
\providecommand{\noteextra}{}

\begin{titlepage}
	\begin{tikzpicture}[remember picture, overlay]
		\begin{scope}[shift={(current page.south west)}]

		% ============================================================
		%  1. 高空→近雾的竖向渐变
		% ============================================================
		\shade[top color=wfSky, bottom color=wfSky2] (0,0) rectangle (21,29.7);

		% ============================================================
		%  2. 雾海中若隐若现的远峰
		% ============================================================
		\fill[wfFar, opacity=0.55]
			(0,16.4) -- (2.4,19.0) -- (3.8,17.6) -- (6.0,20.2) -- (8.4,18.0)
			-- (11.2,20.6) -- (13.8,17.8) -- (16.4,19.8) -- (18.6,17.4)
			-- (21,19.0) -- (21,13.0) -- (0,13.0) -- cycle;
		\fill[wfFar, opacity=0.35]
			(0,13.6) -- (3.0,15.6) -- (5.4,13.8) -- (8.2,16.0) -- (11.6,13.6)
			-- (14.8,15.8) -- (18.0,13.4) -- (21,15.2) -- (21,10.5) -- (0,10.5) -- cycle;

		% ============================================================
		%  3. 翻涌的雾海（横向白色雾带，层层叠压）
		% ============================================================
		\foreach \y/\h/\o in {12.6/1.5/0.85, 10.4/1.3/0.75, 8.2/1.2/0.70,
		                      6.0/1.0/0.60, 4.0/0.9/0.50, 2.2/0.8/0.42}{
			\fill[wfFog, opacity=\o]
				plot[domain=0:21, samples=60, smooth]
					(\x, {\y + 0.30*sin(deg(1.1*\x + 3*\y)) + 0.16*sin(deg(2.3*\x))})
				-- plot[domain=21:0, samples=60, smooth]
					(\x, {\y - \h + 0.26*sin(deg(0.9*\x + 2*\y))})
				-- cycle;
		}

		% ============================================================
		%  4. 中景岩峰（灰，半没于雾中）
		% ============================================================
		\fill[wfRock2, opacity=0.90]
			(0,9.6) -- (2.2,11.2) -- (3.6,10.0) -- (5.2,11.8) -- (6.6,10.2)
			-- (8.0,11.6) -- (9.6,10.0) -- (11.0,11.4) -- (12.6,9.8)
			-- (14.2,11.2) -- (16.0,9.6) -- (18.0,10.8) -- (21,9.4)
			-- (21,0) -- (0,0) -- cycle;

		% ============================================================
		%  5. 前景主峰（深灰，自左下向右上斜出）
		% ============================================================
		\fill[wfRock]
			(0,0) -- (0,7.4) -- (1.8,8.6) -- (3.2,7.6) -- (4.6,9.2)
			-- (6.2,8.2) -- (7.4,9.6) -- (8.6,8.8) -- (9.6,9.7)
			-- (10.6,8.6) -- (11.8,7.0) -- (13.4,6.2) -- (15.6,4.6)
			-- (18.2,3.0) -- (21,1.6) -- (21,0) -- cycle;
		% 岩面裂纹
		\foreach \a/\b/\c/\d in {2.4/8.2/3.4/5.6, 5.2/8.8/6.0/6.2,
		                         7.8/9.2/8.4/6.8, 10.2/8.4/11.2/5.4}{
			\draw[wfFog, opacity=0.18, line width=1.2pt] (\a,\b) -- (\c,\d);
		}

		% ============================================================
		%  6. 漫游者（背影：大衣、头、手杖）
		% ============================================================
		\fill[wfFigure]
			(9.60,13.45)
			.. controls (8.90,13.20) and (8.60,12.20) .. (8.50,11.40)
			.. controls (8.40,10.40) and (8.15,10.00) .. (8.30,9.72)
			.. controls (9.00,9.50) and (10.30,9.50) .. (11.00,9.82)
			.. controls (11.15,10.20) and (10.80,11.20) .. (10.70,12.20)
			.. controls (10.60,13.00) and (10.30,13.40) .. (9.60,13.45) -- cycle;
		\fill[wfFigure] (9.62,14.05) circle[radius=0.52];
		\fill[wfFigure] (9.62,14.42) .. controls (9.30,14.44) and (9.22,14.10) .. (9.30,13.86)
			.. controls (9.60,13.74) and (9.86,13.86) .. (9.62,14.42) -- cycle;
		\draw[wfFigure, line width=2.2pt, line cap=round] (10.85,12.55) -- (11.95,9.45);
		\draw[wfFigure, line width=1.6pt, line cap=round]
			(8.62,11.60) .. controls (8.95,11.75) and (9.25,11.80) .. (9.50,11.78);

		% ============================================================
		%  7. 标题衬底（一层薄雾，保证文字可读）
		% ============================================================
		\fill[wfFog, opacity=0.45]
			(2.6,16.9) -- (17.4,17.3) -- (18.0,22.6) -- (16.4,23.6) -- (2.2,22.9) -- cycle;

		% ============================================================
		%  8. 细边框
		% ============================================================
		\draw[wfInk, line width=0.7pt, opacity=0.5]
			([shift={(0.85,0.85)}]current page.south west) rectangle
			([shift={(-0.85,-0.85)}]current page.north east);

		% ============================================================
		%  9. 顶部眉题（课程英文小字 + 自适应短线）
		% ============================================================
		\node[anchor=north, text=wfInk, align=center]
			(wfEyebrow) at ([yshift=-3.4cm]current page.north) {%
			{\sffamily\fontsize{12}{15}\selectfont \uppercase{\notecourseen}}
		};
		\draw[wfInk, line width=1.0pt] ($(wfEyebrow.west)+(-0.25cm,0)$) -- ++(-1.1cm,0);
		\draw[wfInk, line width=1.0pt] ($(wfEyebrow.east)+(0.25cm,0)$) -- ++(1.1cm,0);

		% ============================================================
		% 10. 主标题（居中，使用 \notename）
		% ============================================================
		\node[anchor=center, align=center, text=wfInk]
			at ([yshift=-8.3cm]current page.north) {%
			{\fontsize{38}{48}\sffamily\bfseries \notename}
		};

		% ============================================================
		% 11. 菱形 + 双细线
		% ============================================================
		\coordinate (C) at ([yshift=-11.6cm]current page.north);
		\draw[wfInk, line width=1.1pt] (C) ++(-2.2cm,0) -- ++(1.55cm,0);
		\draw[wfInk, line width=1.1pt] (C) ++(0.65cm,0) -- ++(1.55cm,0);
		\node[fill=wfInk, minimum size=0.24cm, inner sep=0pt, rotate=45] at (C) {};

		% ============================================================
		% 12. 副标题（使用 \notesubtitle）
		% ============================================================
		\node[anchor=north, text=wfInk, align=center]
			at ([yshift=-12.4cm]current page.north) {
			{\fontsize{15}{20}\sffamily \notesubtitle}
		};

		% ============================================================
		% 13. 底部作者、日期、附加信息
		% ============================================================
		\coordinate (B) at ([yshift=2.7cm]current page.south);
		\draw[wfFog, line width=0.8pt, opacity=0.85] (B) ++(-1.1cm,0.95cm) -- ++(2.2cm,0);
		\node[anchor=south, text=wfFog, align=center] at (B) {
			\begin{tabular}{c}
				{\fontsize{19}{24}\sffamily\bfseries \noteauthor} \\[0.35cm]
				{\large \sffamily \today}
				\ifstrempty{\noteextra}{}{\\[0.3cm]{\normalsize\sffamily \noteextra}}
			\end{tabular}
		};

		% ============================================================
		% 14. 右下角落款（画作信息）
		% ============================================================
		\node[anchor=south east, text=wfFog, opacity=0.8, font=\footnotesize\itshape]
			at ([shift={(-1.35,1.25)}]current page.south east)
			{Caspar David Friedrich, Wanderer above the Sea of Fog, 1818};

		\end{scope}
	\end{tikzpicture}
\end{titlepage}

%
%  配色：wf* 前缀（弗里德里希取色）。本文件用 \providecolor 兜底，
%        单独复制到任何笔记本也能编译；想整体换调改这里即可。
%
%  注意：整幅画面绘于 \begin{scope}[shift={(current page.south west)}] 内，
%        所有坐标以「页面左下角为原点、单位 cm」书写；脱离该 scope
%        直接写 (x,y)，图形会落到页面外而不可见。
%
%  可用字段（在主文件导言区用 \newcommand 预定义即可覆盖缺省值）：
%      \notename     主标题（必填，主模板已定义）
%      \noteauthor   作者（必填，主模板已定义）
%      \notecourseen 眉题英文（缺省 Course Notes）
%      \notesubtitle 副标题（缺省「学习笔记」）
%      \noteextra    底部附加信息（缺省不显示）
% ============================================================

% ---- 弗里德里希《雾海上的漫游者》取色（冷灰蓝 + 墨黑剪影） ----
\providecolor{wfSky}{RGB}{188,198,210}      % 高空：灰蓝（上）
\providecolor{wfSky2}{RGB}{214,220,228}     % 近雾：浅灰蓝
\providecolor{wfFog}{RGB}{240,242,244}      % 雾海：白
\providecolor{wfFar}{RGB}{146,158,172}      % 远峰：淡灰蓝
\providecolor{wfRock2}{RGB}{104,106,112}    % 中景岩：中灰
\providecolor{wfRock}{RGB}{58,60,66}        % 前景岩：深灰
\providecolor{wfFigure}{RGB}{28,30,34}      % 人物：墨黑
\providecolor{wfInk}{RGB}{36,42,50}         % 标题：深蓝灰

% ---- 缺省字段（主文件已定义的同名命令不会被覆盖） ----
\providecommand{\notecourseen}{Course Notes}
\providecommand{\notesubtitle}{学\ 习\ 笔\ 记}
\providecommand{\noteextra}{}

\begin{titlepage}
	\begin{tikzpicture}[remember picture, overlay]
		\begin{scope}[shift={(current page.south west)}]

		% ============================================================
		%  1. 高空→近雾的竖向渐变
		% ============================================================
		\shade[top color=wfSky, bottom color=wfSky2] (0,0) rectangle (21,29.7);

		% ============================================================
		%  2. 雾海中若隐若现的远峰
		% ============================================================
		\fill[wfFar, opacity=0.55]
			(0,16.4) -- (2.4,19.0) -- (3.8,17.6) -- (6.0,20.2) -- (8.4,18.0)
			-- (11.2,20.6) -- (13.8,17.8) -- (16.4,19.8) -- (18.6,17.4)
			-- (21,19.0) -- (21,13.0) -- (0,13.0) -- cycle;
		\fill[wfFar, opacity=0.35]
			(0,13.6) -- (3.0,15.6) -- (5.4,13.8) -- (8.2,16.0) -- (11.6,13.6)
			-- (14.8,15.8) -- (18.0,13.4) -- (21,15.2) -- (21,10.5) -- (0,10.5) -- cycle;

		% ============================================================
		%  3. 翻涌的雾海（横向白色雾带，层层叠压）
		% ============================================================
		\foreach \y/\h/\o in {12.6/1.5/0.85, 10.4/1.3/0.75, 8.2/1.2/0.70,
		                      6.0/1.0/0.60, 4.0/0.9/0.50, 2.2/0.8/0.42}{
			\fill[wfFog, opacity=\o]
				plot[domain=0:21, samples=60, smooth]
					(\x, {\y + 0.30*sin(deg(1.1*\x + 3*\y)) + 0.16*sin(deg(2.3*\x))})
				-- plot[domain=21:0, samples=60, smooth]
					(\x, {\y - \h + 0.26*sin(deg(0.9*\x + 2*\y))})
				-- cycle;
		}

		% ============================================================
		%  4. 中景岩峰（灰，半没于雾中）
		% ============================================================
		\fill[wfRock2, opacity=0.90]
			(0,9.6) -- (2.2,11.2) -- (3.6,10.0) -- (5.2,11.8) -- (6.6,10.2)
			-- (8.0,11.6) -- (9.6,10.0) -- (11.0,11.4) -- (12.6,9.8)
			-- (14.2,11.2) -- (16.0,9.6) -- (18.0,10.8) -- (21,9.4)
			-- (21,0) -- (0,0) -- cycle;

		% ============================================================
		%  5. 前景主峰（深灰，自左下向右上斜出）
		% ============================================================
		\fill[wfRock]
			(0,0) -- (0,7.4) -- (1.8,8.6) -- (3.2,7.6) -- (4.6,9.2)
			-- (6.2,8.2) -- (7.4,9.6) -- (8.6,8.8) -- (9.6,9.7)
			-- (10.6,8.6) -- (11.8,7.0) -- (13.4,6.2) -- (15.6,4.6)
			-- (18.2,3.0) -- (21,1.6) -- (21,0) -- cycle;
		% 岩面裂纹
		\foreach \a/\b/\c/\d in {2.4/8.2/3.4/5.6, 5.2/8.8/6.0/6.2,
		                         7.8/9.2/8.4/6.8, 10.2/8.4/11.2/5.4}{
			\draw[wfFog, opacity=0.18, line width=1.2pt] (\a,\b) -- (\c,\d);
		}

		% ============================================================
		%  6. 漫游者（背影：大衣、头、手杖）
		% ============================================================
		\fill[wfFigure]
			(9.60,13.45)
			.. controls (8.90,13.20) and (8.60,12.20) .. (8.50,11.40)
			.. controls (8.40,10.40) and (8.15,10.00) .. (8.30,9.72)
			.. controls (9.00,9.50) and (10.30,9.50) .. (11.00,9.82)
			.. controls (11.15,10.20) and (10.80,11.20) .. (10.70,12.20)
			.. controls (10.60,13.00) and (10.30,13.40) .. (9.60,13.45) -- cycle;
		\fill[wfFigure] (9.62,14.05) circle[radius=0.52];
		\fill[wfFigure] (9.62,14.42) .. controls (9.30,14.44) and (9.22,14.10) .. (9.30,13.86)
			.. controls (9.60,13.74) and (9.86,13.86) .. (9.62,14.42) -- cycle;
		\draw[wfFigure, line width=2.2pt, line cap=round] (10.85,12.55) -- (11.95,9.45);
		\draw[wfFigure, line width=1.6pt, line cap=round]
			(8.62,11.60) .. controls (8.95,11.75) and (9.25,11.80) .. (9.50,11.78);

		% ============================================================
		%  7. 标题衬底（一层薄雾，保证文字可读）
		% ============================================================
		\fill[wfFog, opacity=0.45]
			(2.6,16.9) -- (17.4,17.3) -- (18.0,22.6) -- (16.4,23.6) -- (2.2,22.9) -- cycle;

		% ============================================================
		%  8. 细边框
		% ============================================================
		\draw[wfInk, line width=0.7pt, opacity=0.5]
			([shift={(0.85,0.85)}]current page.south west) rectangle
			([shift={(-0.85,-0.85)}]current page.north east);

		% ============================================================
		%  9. 顶部眉题（课程英文小字 + 自适应短线）
		% ============================================================
		\node[anchor=north, text=wfInk, align=center]
			(wfEyebrow) at ([yshift=-3.4cm]current page.north) {%
			{\sffamily\fontsize{12}{15}\selectfont \uppercase{\notecourseen}}
		};
		\draw[wfInk, line width=1.0pt] ($(wfEyebrow.west)+(-0.25cm,0)$) -- ++(-1.1cm,0);
		\draw[wfInk, line width=1.0pt] ($(wfEyebrow.east)+(0.25cm,0)$) -- ++(1.1cm,0);

		% ============================================================
		% 10. 主标题（居中，使用 \notename）
		% ============================================================
		\node[anchor=center, align=center, text=wfInk]
			at ([yshift=-8.3cm]current page.north) {%
			{\fontsize{38}{48}\sffamily\bfseries \notename}
		};

		% ============================================================
		% 11. 菱形 + 双细线
		% ============================================================
		\coordinate (C) at ([yshift=-11.6cm]current page.north);
		\draw[wfInk, line width=1.1pt] (C) ++(-2.2cm,0) -- ++(1.55cm,0);
		\draw[wfInk, line width=1.1pt] (C) ++(0.65cm,0) -- ++(1.55cm,0);
		\node[fill=wfInk, minimum size=0.24cm, inner sep=0pt, rotate=45] at (C) {};

		% ============================================================
		% 12. 副标题（使用 \notesubtitle）
		% ============================================================
		\node[anchor=north, text=wfInk, align=center]
			at ([yshift=-12.4cm]current page.north) {
			{\fontsize{15}{20}\sffamily \notesubtitle}
		};

		% ============================================================
		% 13. 底部作者、日期、附加信息
		% ============================================================
		\coordinate (B) at ([yshift=2.7cm]current page.south);
		\draw[wfFog, line width=0.8pt, opacity=0.85] (B) ++(-1.1cm,0.95cm) -- ++(2.2cm,0);
		\node[anchor=south, text=wfFog, align=center] at (B) {
			\begin{tabular}{c}
				{\fontsize{19}{24}\sffamily\bfseries \noteauthor} \\[0.35cm]
				{\large \sffamily \today}
				\ifstrempty{\noteextra}{}{\\[0.3cm]{\normalsize\sffamily \noteextra}}
			\end{tabular}
		};

		% ============================================================
		% 14. 右下角落款（画作信息）
		% ============================================================
		\node[anchor=south east, text=wfFog, opacity=0.8, font=\footnotesize\itshape]
			at ([shift={(-1.35,1.25)}]current page.south east)
			{Caspar David Friedrich, Wanderer above the Sea of Fog, 1818};

		\end{scope}
	\end{tikzpicture}
\end{titlepage}

%
%  配色：wf* 前缀（弗里德里希取色）。本文件用 \providecolor 兜底，
%        单独复制到任何笔记本也能编译；想整体换调改这里即可。
%
%  注意：整幅画面绘于 \begin{scope}[shift={(current page.south west)}] 内，
%        所有坐标以「页面左下角为原点、单位 cm」书写；脱离该 scope
%        直接写 (x,y)，图形会落到页面外而不可见。
%
%  可用字段（在主文件导言区用 \newcommand 预定义即可覆盖缺省值）：
%      \notename     主标题（必填，主模板已定义）
%      \noteauthor   作者（必填，主模板已定义）
%      \notecourseen 眉题英文（缺省 Course Notes）
%      \notesubtitle 副标题（缺省「学习笔记」）
%      \noteextra    底部附加信息（缺省不显示）
% ============================================================

% ---- 弗里德里希《雾海上的漫游者》取色（冷灰蓝 + 墨黑剪影） ----
\providecolor{wfSky}{RGB}{188,198,210}      % 高空：灰蓝（上）
\providecolor{wfSky2}{RGB}{214,220,228}     % 近雾：浅灰蓝
\providecolor{wfFog}{RGB}{240,242,244}      % 雾海：白
\providecolor{wfFar}{RGB}{146,158,172}      % 远峰：淡灰蓝
\providecolor{wfRock2}{RGB}{104,106,112}    % 中景岩：中灰
\providecolor{wfRock}{RGB}{58,60,66}        % 前景岩：深灰
\providecolor{wfFigure}{RGB}{28,30,34}      % 人物：墨黑
\providecolor{wfInk}{RGB}{36,42,50}         % 标题：深蓝灰

% ---- 缺省字段（主文件已定义的同名命令不会被覆盖） ----
\providecommand{\notecourseen}{Course Notes}
\providecommand{\notesubtitle}{学\ 习\ 笔\ 记}
\providecommand{\noteextra}{}

\begin{titlepage}
	\begin{tikzpicture}[remember picture, overlay]
		\begin{scope}[shift={(current page.south west)}]

		% ============================================================
		%  1. 高空→近雾的竖向渐变
		% ============================================================
		\shade[top color=wfSky, bottom color=wfSky2] (0,0) rectangle (21,29.7);

		% ============================================================
		%  2. 雾海中若隐若现的远峰
		% ============================================================
		\fill[wfFar, opacity=0.55]
			(0,16.4) -- (2.4,19.0) -- (3.8,17.6) -- (6.0,20.2) -- (8.4,18.0)
			-- (11.2,20.6) -- (13.8,17.8) -- (16.4,19.8) -- (18.6,17.4)
			-- (21,19.0) -- (21,13.0) -- (0,13.0) -- cycle;
		\fill[wfFar, opacity=0.35]
			(0,13.6) -- (3.0,15.6) -- (5.4,13.8) -- (8.2,16.0) -- (11.6,13.6)
			-- (14.8,15.8) -- (18.0,13.4) -- (21,15.2) -- (21,10.5) -- (0,10.5) -- cycle;

		% ============================================================
		%  3. 翻涌的雾海（横向白色雾带，层层叠压）
		% ============================================================
		\foreach \y/\h/\o in {12.6/1.5/0.85, 10.4/1.3/0.75, 8.2/1.2/0.70,
		                      6.0/1.0/0.60, 4.0/0.9/0.50, 2.2/0.8/0.42}{
			\fill[wfFog, opacity=\o]
				plot[domain=0:21, samples=60, smooth]
					(\x, {\y + 0.30*sin(deg(1.1*\x + 3*\y)) + 0.16*sin(deg(2.3*\x))})
				-- plot[domain=21:0, samples=60, smooth]
					(\x, {\y - \h + 0.26*sin(deg(0.9*\x + 2*\y))})
				-- cycle;
		}

		% ============================================================
		%  4. 中景岩峰（灰，半没于雾中）
		% ============================================================
		\fill[wfRock2, opacity=0.90]
			(0,9.6) -- (2.2,11.2) -- (3.6,10.0) -- (5.2,11.8) -- (6.6,10.2)
			-- (8.0,11.6) -- (9.6,10.0) -- (11.0,11.4) -- (12.6,9.8)
			-- (14.2,11.2) -- (16.0,9.6) -- (18.0,10.8) -- (21,9.4)
			-- (21,0) -- (0,0) -- cycle;

		% ============================================================
		%  5. 前景主峰（深灰，自左下向右上斜出）
		% ============================================================
		\fill[wfRock]
			(0,0) -- (0,7.4) -- (1.8,8.6) -- (3.2,7.6) -- (4.6,9.2)
			-- (6.2,8.2) -- (7.4,9.6) -- (8.6,8.8) -- (9.6,9.7)
			-- (10.6,8.6) -- (11.8,7.0) -- (13.4,6.2) -- (15.6,4.6)
			-- (18.2,3.0) -- (21,1.6) -- (21,0) -- cycle;
		% 岩面裂纹
		\foreach \a/\b/\c/\d in {2.4/8.2/3.4/5.6, 5.2/8.8/6.0/6.2,
		                         7.8/9.2/8.4/6.8, 10.2/8.4/11.2/5.4}{
			\draw[wfFog, opacity=0.18, line width=1.2pt] (\a,\b) -- (\c,\d);
		}

		% ============================================================
		%  6. 漫游者（背影：大衣、头、手杖）
		% ============================================================
		\fill[wfFigure]
			(9.60,13.45)
			.. controls (8.90,13.20) and (8.60,12.20) .. (8.50,11.40)
			.. controls (8.40,10.40) and (8.15,10.00) .. (8.30,9.72)
			.. controls (9.00,9.50) and (10.30,9.50) .. (11.00,9.82)
			.. controls (11.15,10.20) and (10.80,11.20) .. (10.70,12.20)
			.. controls (10.60,13.00) and (10.30,13.40) .. (9.60,13.45) -- cycle;
		\fill[wfFigure] (9.62,14.05) circle[radius=0.52];
		\fill[wfFigure] (9.62,14.42) .. controls (9.30,14.44) and (9.22,14.10) .. (9.30,13.86)
			.. controls (9.60,13.74) and (9.86,13.86) .. (9.62,14.42) -- cycle;
		\draw[wfFigure, line width=2.2pt, line cap=round] (10.85,12.55) -- (11.95,9.45);
		\draw[wfFigure, line width=1.6pt, line cap=round]
			(8.62,11.60) .. controls (8.95,11.75) and (9.25,11.80) .. (9.50,11.78);

		% ============================================================
		%  7. 标题衬底（一层薄雾，保证文字可读）
		% ============================================================
		\fill[wfFog, opacity=0.45]
			(2.6,16.9) -- (17.4,17.3) -- (18.0,22.6) -- (16.4,23.6) -- (2.2,22.9) -- cycle;

		% ============================================================
		%  8. 细边框
		% ============================================================
		\draw[wfInk, line width=0.7pt, opacity=0.5]
			([shift={(0.85,0.85)}]current page.south west) rectangle
			([shift={(-0.85,-0.85)}]current page.north east);

		% ============================================================
		%  9. 顶部眉题（课程英文小字 + 自适应短线）
		% ============================================================
		\node[anchor=north, text=wfInk, align=center]
			(wfEyebrow) at ([yshift=-3.4cm]current page.north) {%
			{\sffamily\fontsize{12}{15}\selectfont \uppercase{\notecourseen}}
		};
		\draw[wfInk, line width=1.0pt] ($(wfEyebrow.west)+(-0.25cm,0)$) -- ++(-1.1cm,0);
		\draw[wfInk, line width=1.0pt] ($(wfEyebrow.east)+(0.25cm,0)$) -- ++(1.1cm,0);

		% ============================================================
		% 10. 主标题（居中，使用 \notename）
		% ============================================================
		\node[anchor=center, align=center, text=wfInk]
			at ([yshift=-8.3cm]current page.north) {%
			{\fontsize{38}{48}\sffamily\bfseries \notename}
		};

		% ============================================================
		% 11. 菱形 + 双细线
		% ============================================================
		\coordinate (C) at ([yshift=-11.6cm]current page.north);
		\draw[wfInk, line width=1.1pt] (C) ++(-2.2cm,0) -- ++(1.55cm,0);
		\draw[wfInk, line width=1.1pt] (C) ++(0.65cm,0) -- ++(1.55cm,0);
		\node[fill=wfInk, minimum size=0.24cm, inner sep=0pt, rotate=45] at (C) {};

		% ============================================================
		% 12. 副标题（使用 \notesubtitle）
		% ============================================================
		\node[anchor=north, text=wfInk, align=center]
			at ([yshift=-12.4cm]current page.north) {
			{\fontsize{15}{20}\sffamily \notesubtitle}
		};

		% ============================================================
		% 13. 底部作者、日期、附加信息
		% ============================================================
		\coordinate (B) at ([yshift=2.7cm]current page.south);
		\draw[wfFog, line width=0.8pt, opacity=0.85] (B) ++(-1.1cm,0.95cm) -- ++(2.2cm,0);
		\node[anchor=south, text=wfFog, align=center] at (B) {
			\begin{tabular}{c}
				{\fontsize{19}{24}\sffamily\bfseries \noteauthor} \\[0.35cm]
				{\large \sffamily \today}
				\ifstrempty{\noteextra}{}{\\[0.3cm]{\normalsize\sffamily \noteextra}}
			\end{tabular}
		};

		% ============================================================
		% 14. 右下角落款（画作信息）
		% ============================================================
		\node[anchor=south east, text=wfFog, opacity=0.8, font=\footnotesize\itshape]
			at ([shift={(-1.35,1.25)}]current page.south east)
			{Caspar David Friedrich, Wanderer above the Sea of Fog, 1818};

		\end{scope}
	\end{tikzpicture}
\end{titlepage}

%
%  配色：wf* 前缀（弗里德里希取色）。本文件用 \providecolor 兜底，
%        单独复制到任何笔记本也能编译；想整体换调改这里即可。
%
%  注意：整幅画面绘于 \begin{scope}[shift={(current page.south west)}] 内，
%        所有坐标以「页面左下角为原点、单位 cm」书写；脱离该 scope
%        直接写 (x,y)，图形会落到页面外而不可见。
%
%  可用字段（在主文件导言区用 \newcommand 预定义即可覆盖缺省值）：
%      \notename     主标题（必填，主模板已定义）
%      \noteauthor   作者（必填，主模板已定义）
%      \notecourseen 眉题英文（缺省 Course Notes）
%      \notesubtitle 副标题（缺省「学习笔记」）
%      \noteextra    底部附加信息（缺省不显示）
% ============================================================

% ---- 弗里德里希《雾海上的漫游者》取色（冷灰蓝 + 墨黑剪影） ----
\providecolor{wfSky}{RGB}{188,198,210}      % 高空：灰蓝（上）
\providecolor{wfSky2}{RGB}{214,220,228}     % 近雾：浅灰蓝
\providecolor{wfFog}{RGB}{240,242,244}      % 雾海：白
\providecolor{wfFar}{RGB}{146,158,172}      % 远峰：淡灰蓝
\providecolor{wfRock2}{RGB}{104,106,112}    % 中景岩：中灰
\providecolor{wfRock}{RGB}{58,60,66}        % 前景岩：深灰
\providecolor{wfFigure}{RGB}{28,30,34}      % 人物：墨黑
\providecolor{wfInk}{RGB}{36,42,50}         % 标题：深蓝灰

% ---- 缺省字段（主文件已定义的同名命令不会被覆盖） ----
\providecommand{\notecourseen}{Course Notes}
\providecommand{\notesubtitle}{学\ 习\ 笔\ 记}
\providecommand{\noteextra}{}

\begin{titlepage}
	\begin{tikzpicture}[remember picture, overlay]
		\begin{scope}[shift={(current page.south west)}]

		% ============================================================
		%  1. 高空→近雾的竖向渐变
		% ============================================================
		\shade[top color=wfSky, bottom color=wfSky2] (0,0) rectangle (21,29.7);

		% ============================================================
		%  2. 雾海中若隐若现的远峰
		% ============================================================
		\fill[wfFar, opacity=0.55]
			(0,16.4) -- (2.4,19.0) -- (3.8,17.6) -- (6.0,20.2) -- (8.4,18.0)
			-- (11.2,20.6) -- (13.8,17.8) -- (16.4,19.8) -- (18.6,17.4)
			-- (21,19.0) -- (21,13.0) -- (0,13.0) -- cycle;
		\fill[wfFar, opacity=0.35]
			(0,13.6) -- (3.0,15.6) -- (5.4,13.8) -- (8.2,16.0) -- (11.6,13.6)
			-- (14.8,15.8) -- (18.0,13.4) -- (21,15.2) -- (21,10.5) -- (0,10.5) -- cycle;

		% ============================================================
		%  3. 翻涌的雾海（横向白色雾带，层层叠压）
		% ============================================================
		\foreach \y/\h/\o in {12.6/1.5/0.85, 10.4/1.3/0.75, 8.2/1.2/0.70,
		                      6.0/1.0/0.60, 4.0/0.9/0.50, 2.2/0.8/0.42}{
			\fill[wfFog, opacity=\o]
				plot[domain=0:21, samples=60, smooth]
					(\x, {\y + 0.30*sin(deg(1.1*\x + 3*\y)) + 0.16*sin(deg(2.3*\x))})
				-- plot[domain=21:0, samples=60, smooth]
					(\x, {\y - \h + 0.26*sin(deg(0.9*\x + 2*\y))})
				-- cycle;
		}

		% ============================================================
		%  4. 中景岩峰（灰，半没于雾中）
		% ============================================================
		\fill[wfRock2, opacity=0.90]
			(0,9.6) -- (2.2,11.2) -- (3.6,10.0) -- (5.2,11.8) -- (6.6,10.2)
			-- (8.0,11.6) -- (9.6,10.0) -- (11.0,11.4) -- (12.6,9.8)
			-- (14.2,11.2) -- (16.0,9.6) -- (18.0,10.8) -- (21,9.4)
			-- (21,0) -- (0,0) -- cycle;

		% ============================================================
		%  5. 前景主峰（深灰，自左下向右上斜出）
		% ============================================================
		\fill[wfRock]
			(0,0) -- (0,7.4) -- (1.8,8.6) -- (3.2,7.6) -- (4.6,9.2)
			-- (6.2,8.2) -- (7.4,9.6) -- (8.6,8.8) -- (9.6,9.7)
			-- (10.6,8.6) -- (11.8,7.0) -- (13.4,6.2) -- (15.6,4.6)
			-- (18.2,3.0) -- (21,1.6) -- (21,0) -- cycle;
		% 岩面裂纹
		\foreach \a/\b/\c/\d in {2.4/8.2/3.4/5.6, 5.2/8.8/6.0/6.2,
		                         7.8/9.2/8.4/6.8, 10.2/8.4/11.2/5.4}{
			\draw[wfFog, opacity=0.18, line width=1.2pt] (\a,\b) -- (\c,\d);
		}

		% ============================================================
		%  6. 漫游者（背影：大衣、头、手杖）
		% ============================================================
		\fill[wfFigure]
			(9.60,13.45)
			.. controls (8.90,13.20) and (8.60,12.20) .. (8.50,11.40)
			.. controls (8.40,10.40) and (8.15,10.00) .. (8.30,9.72)
			.. controls (9.00,9.50) and (10.30,9.50) .. (11.00,9.82)
			.. controls (11.15,10.20) and (10.80,11.20) .. (10.70,12.20)
			.. controls (10.60,13.00) and (10.30,13.40) .. (9.60,13.45) -- cycle;
		\fill[wfFigure] (9.62,14.05) circle[radius=0.52];
		\fill[wfFigure] (9.62,14.42) .. controls (9.30,14.44) and (9.22,14.10) .. (9.30,13.86)
			.. controls (9.60,13.74) and (9.86,13.86) .. (9.62,14.42) -- cycle;
		\draw[wfFigure, line width=2.2pt, line cap=round] (10.85,12.55) -- (11.95,9.45);
		\draw[wfFigure, line width=1.6pt, line cap=round]
			(8.62,11.60) .. controls (8.95,11.75) and (9.25,11.80) .. (9.50,11.78);

		% ============================================================
		%  7. 标题衬底（一层薄雾，保证文字可读）
		% ============================================================
		\fill[wfFog, opacity=0.45]
			(2.6,16.9) -- (17.4,17.3) -- (18.0,22.6) -- (16.4,23.6) -- (2.2,22.9) -- cycle;

		% ============================================================
		%  8. 细边框
		% ============================================================
		\draw[wfInk, line width=0.7pt, opacity=0.5]
			([shift={(0.85,0.85)}]current page.south west) rectangle
			([shift={(-0.85,-0.85)}]current page.north east);

		% ============================================================
		%  9. 顶部眉题（课程英文小字 + 自适应短线）
		% ============================================================
		\node[anchor=north, text=wfInk, align=center]
			(wfEyebrow) at ([yshift=-3.4cm]current page.north) {%
			{\sffamily\fontsize{12}{15}\selectfont \uppercase{\notecourseen}}
		};
		\draw[wfInk, line width=1.0pt] ($(wfEyebrow.west)+(-0.25cm,0)$) -- ++(-1.1cm,0);
		\draw[wfInk, line width=1.0pt] ($(wfEyebrow.east)+(0.25cm,0)$) -- ++(1.1cm,0);

		% ============================================================
		% 10. 主标题（居中，使用 \notename）
		% ============================================================
		\node[anchor=center, align=center, text=wfInk]
			at ([yshift=-8.3cm]current page.north) {%
			{\fontsize{38}{48}\sffamily\bfseries \notename}
		};

		% ============================================================
		% 11. 菱形 + 双细线
		% ============================================================
		\coordinate (C) at ([yshift=-11.6cm]current page.north);
		\draw[wfInk, line width=1.1pt] (C) ++(-2.2cm,0) -- ++(1.55cm,0);
		\draw[wfInk, line width=1.1pt] (C) ++(0.65cm,0) -- ++(1.55cm,0);
		\node[fill=wfInk, minimum size=0.24cm, inner sep=0pt, rotate=45] at (C) {};

		% ============================================================
		% 12. 副标题（使用 \notesubtitle）
		% ============================================================
		\node[anchor=north, text=wfInk, align=center]
			at ([yshift=-12.4cm]current page.north) {
			{\fontsize{15}{20}\sffamily \notesubtitle}
		};

		% ============================================================
		% 13. 底部作者、日期、附加信息
		% ============================================================
		\coordinate (B) at ([yshift=2.7cm]current page.south);
		\draw[wfFog, line width=0.8pt, opacity=0.85] (B) ++(-1.1cm,0.95cm) -- ++(2.2cm,0);
		\node[anchor=south, text=wfFog, align=center] at (B) {
			\begin{tabular}{c}
				{\fontsize{19}{24}\sffamily\bfseries \noteauthor} \\[0.35cm]
				{\large \sffamily \today}
				\ifstrempty{\noteextra}{}{\\[0.3cm]{\normalsize\sffamily \noteextra}}
			\end{tabular}
		};

		% ============================================================
		% 14. 右下角落款（画作信息）
		% ============================================================
		\node[anchor=south east, text=wfFog, opacity=0.8, font=\footnotesize\itshape]
			at ([shift={(-1.35,1.25)}]current page.south east)
			{Caspar David Friedrich, Wanderer above the Sea of Fog, 1818};

		\end{scope}
	\end{tikzpicture}
\end{titlepage}
